\documentclass[11pt,a4paper]{article}
\usepackage[T1]{fontenc}
\usepackage[utf8]{inputenc}
\usepackage{lmodern}
\usepackage{amsmath,amssymb}
\usepackage[margin=25mm]{geometry}
\usepackage[hidelinks]{hyperref}
\hypersetup{pdftitle={From the continuum to discrete substrates: 't Hooft and other turns},pdfauthor={navstokgap research working notes}}
\newcommand{\tightlist}{\setlength{\itemsep}{0pt}\setlength{\parskip}{0pt}}
\setlength{\emergencystretch}{3em}
\setcounter{secnumdepth}{0}
\begin{document}
\hypertarget{from-the-continuum-to-discrete-substrates-t-hooft-and-other-turns}{%
\section{From the continuum to discrete substrates: 't Hooft and other
turns}\label{from-the-continuum-to-discrete-substrates-t-hooft-and-other-turns}}

\textbf{Conversation record, 2026-09-28 (user and Claude, with a
from-memory answer by Claude Fable).} All references below are quoted
from memory and have not been verified; they are leads for the prior-art
section of the fifth-postulate paper, and no result here is claimed.

Conversation with the user, 2026-09-28 night (not yet in any repo note;
references from memory, unverified).

\textbf{'t Hooft} (Fable's memory, high/medium confidence):
renormalizability and beta\textless0 (Marseille 1972); only first two
beta coefficients scheme-independent (limits ``2 b0 per log'' beyond one
loop); renormalons (Erice late 1970s) = the atlas exponential ladder,
threshold 2 b0 c = 4 is the gluon-condensate renormalon; instantons
e\^{}\{-8 pi\textsuperscript{2/g}2\}, U(1) problem; 't Hooft loop B(C)
(NPB 138, 1978) = our mid-edge centre flip; twisted b.c. and e/m flux
sectors (NPB 153, 1979), fractional charge 1/N mod 1 (1/3 for SU(3),
fits the trisection result); the round-12 non-based-cycle Aharonov-Bohm
phase is the small-angle part of his flux-sector data (round-14 prompt
now asks Astra to state the covariant reference sector by sector);
abelian projection 1981; large N 1974. Cellular-automaton interpretation
(1988 on, book 2016): hbar enters as a conversion constant tied to the
time grain delta-t; the floor comes from information loss (equivalence
classes give H bounded below). This is the exact opposite of the user's
thesis on the same RG ground: he supplies hbar from below; we ask
whether consistent refinement forces it.

\textbf{Candidate next Newton premise} (proposed, not yet sent to
Astra): take 't Hooft's information-loss/equivalence-class construction
as the physical premise and ask whether it bounds the recorded Galileo
comparison. Contact point: our thermodynamic-records note gave only eta
\textgreater= A0 e\^{}\{-W/kBT\}, no floor. Put this to an Astra
decision round after round 14.

\textbf{User's remarks:} ``as soon as you really understand the RG then
you don't understand any'' (amused); read Wolfram in the eighties
thinking ``how will he reach a continuum limit'', yet the Newton thesis
says the continuum limit is problematic from the start, so ``a bit
bipolar''. Claude's reply: same fault line from two sides; the atlas
takes the limit and asks what survives.

\textbf{Wolfram} (Caltech \textasciitilde1979, QCD phenomenology,
Fox-Wolfram moments; cellular automata 1981-82; NKS 2002; hypergraph
physics 2020). \textbf{Cvitanovic} (Kinoshita's group, QED g-2;
birdtracks from colour factors, book 2008, useful notation for our SU(N)
Brauer-Klimyk sums; chaos via the Feigenbaum-Cvitanovic RG fixed point,
ChaosBook).

\textbf{Other ``Damascus'' actors listed:} divergence-driven: Heisenberg
universal length (1938), Snyder quantized spacetime (1947), Yukawa
elementary domains, Dirac's lifelong rejection of renormalization. RG
masters going deterministic/discrete: T. D. Lee discrete time (1983),
Adler trace dynamics (2004), H. B. Nielsen random dynamics, Feynman
checkerboard and ``Simulating physics with computers'' (1982). Geometry
to combinatorics: Regge, Penrose spin networks, Finkelstein space-time
code (1969), Sorkin causal sets, Ambjorn-Loll, Wheeler ``it from bit'',
von Weizsacker ur-alternatives. \textbf{Closest to the Newton thesis:}
Gisin and Del Santo (\textasciitilde2019-21), finite information makes
classical mechanics indeterministic (denial of the velocitas ultima from
quantum information) but give no hbar; the fifth-postulate paper should
cite and answer them. Proposed paper sentence: field theorists went
discrete to escape infinities; RG masters to explain the quantum;
Gisin/Del Santo to break determinism; the user's thesis keeps the
continuum and shows consistent refinement demands the constant Newton's
limit dropped.

Related: \href{../research/STATE.md}{STATE},
\href{principia-fifth-postulate.md}{fifth postulate},
\href{stochastic-route-velocitas-ultima.md}{stochastic route},
\href{thermodynamic-records-no-floor.md}{thermodynamic records}.

\hypertarget{t-hoofts-automaton-against-zeno-the-atomists-and-newton}{%
\subsection{'t Hooft's automaton against Zeno, the atomists and
Newton}\label{t-hoofts-automaton-against-zeno-the-atomists-and-newton}}

Added 2026-09-29 (Claude), at the user's request. The automaton papers
are identified by DOI (Crossref-verified, metadata only): the
\href{https://doi.org/10.1007/BF01011560}{1988 equivalence relations},
the \href{https://doi.org/10.1088/0264-9381/16/10/316}{1999 dissipative
deterministic system} and the
\href{https://doi.org/10.1007/978-3-319-41285-6}{2016 book}. What
follows is a reading of his position as Fable and Claude remember it,
set against the classics entries of the
\href{planck-gap-paper.md}{Planck paper} §§7 and 9.

\textbf{His commitments.} Reality is a deterministic automaton with a
universal time step \(\delta t\). A reversible automaton is a
permutation of states, so its evolution operator has eigenvalues on the
unit circle and energies defined only modulo \(2\pi\hbar/\delta t\):
\(\hbar\) is the conversion between the step and an energy, and the
spectrum has no ground state. Information loss, many states merging into
one equivalence class, is what he invokes to obtain a Hamiltonian
bounded below. Quantum mechanics is the description of the classes; Bell
correlations are paid for by superdeterminism.

\begin{itemize}
\tightlist
\item
  \textbf{Zeno and Diodorus.} At a tick the automaton holds a
  configuration and motion exists only as the update between ticks:
  Diodorus's ``a thing never is moving, but it has moved'', made into
  dynamics. The Planck paper keeps Zeno's conclusion for local records
  only, below the window \(\tau_{\rm arrow}=8z^2\kappa/(mv^2)\), which
  depends on the body and the confidence. The automaton makes the arrow
  exact at every scale below \(\delta t\) and for every body.
\item
  \textbf{Democritus and the kalam atom.} The automaton's cells and
  ticks are atoms of place and time, the kalam position on time and a
  stepped cone in Democritus's dilemma. The Planck paper denies the
  universal time atom on the evidence of records: its mesh moves with
  \(F\), \(m\) and the confidence, a dynamical resolution. The two are
  compatible: a universal grain far below every mark mesh leaves Theorem
  9 untouched. They differ in status: the grain is ontological, the mesh
  is epistemic.
\item
  \textbf{Epicurus.} Continuity presented to sense, succession below it,
  and a warning against extrapolating continuity downward: this is 't
  Hooft's picture with the Planck time in place of the threshold of
  sense. Epicurus is the closest ancient statement of it, as he is of
  ours; the difference is where the boundary is drawn and whether it is
  derived.
\item
  \textbf{Newton.} The automaton removes the \emph{velocitas ultima} by
  removing the limit: velocities are finite differences and no instant
  carries one. That is a third denial of joint determinacy, beside
  quantum noncommutativity (Theorem A) and stochastic paths (the
  \href{stochastic-route-velocitas-ultima.md}{stochastic route}), and it
  keeps determinism. Its \(\hbar\) is a unit conversion, so the
  automaton alone supplies a grain and no action floor on records: a
  reversible automaton is a discrete Liouville dynamics, and Theorem I
  of the
  \href{necessity-unit-and-indeterminacy.md}{unit-and-indeterminacy
  note} gives classical readouts no floor. The floor, in his
  construction, rides on information loss, and the
  \href{thermodynamic-records-no-floor.md}{thermodynamic-records note}
  finds that loss yields a trade-off \(\eta\ge A_0e^{-W/k_BT}\) and no
  floor.
\end{itemize}

\textbf{The sharp question this leaves}, the candidate Newton premise
above: does coarse-graining by equivalence classes, of the kind that
gives his Hamiltonian a ground state, force a positive floor on the
recorded inertial--parabola comparison, or only a trade-off? The
repository's current results point to the second; a theorem either way
would place the automaton exactly relative to the thesis that
quantization is a consistency condition of the continuum limit.

\hypertarget{later-the-same-night-superdeterminism-and-dedekinds-cut}{%
\subsection{Later the same night: superdeterminism and Dedekind's
cut}\label{later-the-same-night-superdeterminism-and-dedekinds-cut}}

\textbf{Superdeterminism} (user: ``classical mechanics is accidentally
superdeterminist too, and we think it has some mistake somewhere'').
From memory: Bell called the independence of settings and hidden
variables ``free variables'' (exchange with Shimony, Horne and Clauser,
around 1976--77) and named the escape ``super-deterministic'' in a 1985
BBC interview (\emph{The Ghost in the Atom}, 1986); Brans gave a model
in 1988; 't Hooft adopted it in the 2000s; Hossenfelder and Palmer
revived it in 2020. The link the user and Claude found interesting is
developed in the \href{superdeterminism-floor.md}{superdeterminism
note}: for chaotic setting devices, a superdeterministic correlation
must be stored in structure of the initial data finer than any fixed
resolution, so a floor gives it an Ehrenfest-type lifetime and \(h\to0\)
reopens the loophole.

\textbf{Dedekind's cut} (user: ``Dedekind cuts are peculiar, as they
define a point as a segment, really''). Dedekind (1872) defines a real
number by the two segments of rationals it separates: the point is given
by what lies on either side, and fixing it needs infinitely many
comparisons. The intuitionists, and after them Gisin and Del Santo (from
memory: finite-information quantities whose digits are fixed
progressively, around 2019--21), refuse the completed cut: a real is
only ever a nested sequence of segments, never a point. The repository
has a quantitative form of the remark. Locating an instant by nested
cuts spends Newton's cell action additively
(\href{cut-measure-newton.md}{cut-measure note}, Theorem 2), and with a
floor \(\kappa\) per exhibited cut only \(K_\tau/\kappa\) cuts can be
exhibited (Corollary 5): a recorded instant is a segment of length of
order the mark mesh \(\tau_*\). Dedekind's point is the completed limit
of that sequence, which is Newton's side in the scholium closing Book I
(Euclid X against least magnitudes); the recorded point is the
unfinished one.

\textbf{Connes} (user: ``Connes more or less exits the dilemma by having
\(dx=[D,X]\)''). From memory: in a spectral triple the differential of
\(f\) is \([D,f]\), infinitesimals are compact operators, the line
element is \(ds=D^{-1}\), and distance needs no paths,
\(d(p,q)=\sup\{|f(p)-f(q)|:\|[D,f]\|\le1\}\). The completed point is
never formed, and the velocity is an operator relation, as in
Heisenberg's \(\dot x=(i/\hbar)[H,x]\). Two refinements: in the
commutative case \([D,f]\) commutes with every function, so joint
determinacy holds and Newton is recovered; noncommutativity remains the
premise, the fork of Theorem A of the
\href{principia-fifth-postulate.md}{fifth-postulate note}. With
\(p=\hbar D\) the line element \(D^{-1}\) is \(\hbar/p\), the de Broglie
wavelength operator, so \(\hbar\) is again a conversion constant. The
\href{tangent-groupoid-trajectories.md}{tangent-groupoid note} already
holds the gluing of pairs at \(\hbar>0\) (the segment) to tangent
vectors at \(\hbar=0\) (the point). A candidate Newton test: whether the
repeated cutting of a Galileo cell extends continuously to the
\(\hbar=0\) fibre.

\hypertarget{consequence-for-state}{%
\subsection{Consequence for STATE}\label{consequence-for-state}}

None yet. The candidate Newton premise (information loss as in 't
Hooft's equivalence classes) goes to an Astra decision round after round
14.
\end{document}
